\documentclass[11pt]{amsart}

\usepackage[utf8]{inputenc}
\usepackage[T1]{fontenc}
\usepackage[margin=1in]{geometry}
\usepackage{amsmath,amssymb,amsthm,mathrsfs}
\usepackage{booktabs}
\usepackage{hyperref}
\usepackage{cleveref}

\newtheorem{theorem}{Theorem}
\newtheorem{proposition}[theorem]{Proposition}
\newtheorem{lemma}[theorem]{Lemma}
\newtheorem{corollary}[theorem]{Corollary}
\theoremstyle{definition}
\newtheorem{definition}[theorem]{Definition}
\theoremstyle{remark}
\newtheorem{remark}[theorem]{Remark}

\DeclareMathOperator{\Cl}{Cl}
\DeclareMathOperator{\Gal}{Gal}
\DeclareMathOperator{\Tr}{Tr}
\DeclareMathOperator{\Frob}{Frob}
\newcommand{\Q}{\mathbb{Q}}
\newcommand{\Z}{\mathbb{Z}}
\newcommand{\C}{\mathbb{C}}
\newcommand{\OK}{\mathcal{O}_K}
\newcommand{\chiK}{\chi_K}
\newcommand{\frakp}{\mathfrak{p}}
\newcommand{\fraka}{\mathfrak{a}}
\newcommand{\frakf}{\mathfrak{f}}

\title[Newton--Dickson recursion for CM newforms]{A Newton--Dickson recursion for Hecke eigenvalues of
CM newforms over imaginary quadratic fields}

\author[K. R\'emondi\`ere]{K\'evin R\'emondi\`ere\\
\small Independent Researcher, Oloron-Sainte-Marie, France\\
\small ORCID: \href{https://orcid.org/0009-0008-2443-7166}{0009-0008-2443-7166}\\
\small \texttt{kevin.remondiere@gmail.com}}
\address{Independent researcher}
\date{\today}

\begin{document}

\begin{abstract}
Let $K = \Q(\sqrt{D})$ be an imaginary quadratic field of class
number one and fundamental discriminant $D < 0$, and let
$\{f_w\}_{w \ge 3}$ be the family of weight-$w$ $\Q$-rational
CM newforms attached to a fixed Hecke Gr\"o\ss encharakter $\psi$ of
infinity type $(1,0)$ of $K$. We record the elementary observation
that the Hecke eigenvalues at all unramified primes $p$ satisfy the
two-step linear recursion
\[
  a_p(f_w) \;=\; X(p)\, a_p(f_{w-1}) \;-\; p\, a_p(f_{w-2})
  \qquad (w \ge 3),
\]
with seed $X(p) := a_p(f_2) = \alpha_p + \beta_p$ (the trace of
Frobenius on the associated CM elliptic curve $E_K/\Q$), and closed
form $a_p(f_w) = D_{w-1}(X(p), p)$ where $D_n$ is the Dickson
polynomial of the first kind. The result is a direct consequence of
the Hecke--Shimura framework combined with Newton's identity for
power sums of two variables, and is implicit in
\cite{Shimura1971,Hecke1936,Ribet1977}; we extract it as a named
identity and verify it numerically on $57$ split-prime instances
across discriminants $D \in \{-7, -11, -19, -43, -67, -163\}$ and
weights $w \in \{3,5,7,9,11\}$ using PARI/GP. A componentwise
extension to $h_K \ge 2$ is recorded in \Cref{sec:extension}.
\end{abstract}

\maketitle

\section{Introduction}\label{sec:intro}

Let $K = \Q(\sqrt{D})$ be an imaginary quadratic field with $D < 0$
a fundamental discriminant and ring of integers $\OK$. Throughout
this note we assume $h_K = |\Cl(K)| = 1$, i.e.\ $D \in \{-3, -4, -7,
-8, -11, -19, -43, -67, -163\}$, the nine class-number-one imaginary
quadratic discriminants. Let $\chiK = \bigl(\tfrac{D}{\cdot}\bigr)$
denote the Kronecker character of $K$.

By Hecke \cite{Hecke1936} and Shimura
\cite[\S\S3.4, 7.6]{Shimura1971}, every algebraic Hecke
Gr\"o\ss encharakter $\psi$ of $K$ of infinity type $(w-1, 0)$ and
finite-order central character gives rise to a holomorphic CM
newform
\begin{equation}\label{eq:theta}
  f_w(\tau) \;=\; \sum_{\fraka \subset \OK} \psi(\fraka)\,
  q^{N(\fraka)}, \qquad q = e^{2\pi i \tau},
\end{equation}
of weight $w$ on $\Gamma_0(|D|N(\frakf))$ with character
$\chiK \cdot \chi'$, where $\frakf$ is the conductor of $\psi$ and
$\chi'$ is a finite-order character. Fixing a base character
$\psi_2$ of infinity type $(1, 0)$ (which exists for $h_K = 1$ and
is unique up to its dual; cf.\ \cite{Deuring1953}), the family
$\psi_w := \psi_2^{w-1}$ for $w \ge 2$ yields a tower of CM
newforms $\{f_w\}_{w \ge 2}$ with the same conductor.

For a rational prime $p \nmid |D| N(\frakf)$ split in $\OK$ as $p\OK
= \frakp \bar\frakp$, write $\alpha_p := \psi_2(\frakp)$ and
$\beta_p := \psi_2(\bar\frakp)$, so that $\alpha_p \beta_p =
\psi_2((p)) = p$ in the corpus normalisation. The Hecke eigenvalue
of $f_w$ at $p$ is
\begin{equation}\label{eq:apf-split}
  a_p(f_w) \;=\; \psi_w(\frakp) + \psi_w(\bar\frakp)
  \;=\; \alpha_p^{w-1} + \beta_p^{w-1}.
\end{equation}
At inert primes $p$ the eigenvalue $a_p(f_w)$ vanishes by the
Hecke-character framework, and at the unique ramified prime
$p = |D|$ (for fundamental $D$ with $\frakf = (\sqrt{D})$) the
canonical $\Q$-rational form has $a_p(f_w) = 0$ as well; see
\Cref{sec:inert-ramified}.

\subsection{Statement of the main result}

\begin{theorem}\label{thm:main}
Let $K = \Q(\sqrt{D})$ be imaginary quadratic of class number one,
let $\psi_2$ be the canonical Hecke character of infinity type
$(1, 0)$ of $K$, and let $\{f_w\}_{w \ge 2}$ be the associated
tower of CM newforms with $\psi_w = \psi_2^{w-1}$. Set
$X(p) := a_p(f_2) = \alpha_p + \beta_p$ for $p$ split in $\OK$, and
$X(p) := 0$ for $p$ inert or ramified.

Then for every prime $p \nmid |D| N(\frakf)$, the Hecke eigenvalues
$\{a_p(f_w)\}_{w \ge 1}$ satisfy the linear recursion
\begin{equation}\label{eq:NCseq}
  a_p(f_w) \;=\; X(p)\, a_p(f_{w-1}) \;-\; p\, a_p(f_{w-2})
  \qquad (w \ge 3),
\end{equation}
with initial values $a_p(f_1) = 2$ and $a_p(f_2) = X(p)$.
\end{theorem}

We refer to \cref{eq:NCseq} as the \emph{Newton--Dickson recursion},
in honour of Newton's identity for power sums and Dickson's
\cite{Dickson1897} polynomial reformulation. The equivalent closed
form is

\begin{corollary}\label{cor:dickson}
Under the hypotheses of \cref{thm:main},
\[
  a_p(f_w) \;=\; D_{w-1}\bigl(X(p), p\bigr),
\]
where $D_n(X, p) := \sum_{k=0}^{\lfloor n/2 \rfloor}
\tfrac{n}{n-k} \binom{n-k}{k} (-p)^k X^{n-2k}$ is the Dickson
polynomial of the first kind. Equivalently, $a_p(f_w) =
V_{w-1}(X(p), p)$ where $V_n(P, Q)$ is the Lucas companion sequence
$V_0 = 2$, $V_1 = P$, $V_n = P V_{n-1} - Q V_{n-2}$.
\end{corollary}

The recursion \cref{eq:NCseq} is implicit in \cite{Hecke1936,
Shimura1971, Ribet1977}; the Dickson--Chebyshev reformulation is
folklore in symmetric-function theory \cite[\S I.2]{Macdonald1995}.
Our modest contribution is to extract \cref{thm:main} as a named
identity, give a self-contained three-line proof, and verify it
numerically on the LMFDB tower
$67.3.b.a \to 67.5.b.a \to 67.7.b.a \to 67.9.b.a \to 67.11.b.a$ at
$D = -67$ together with cross-discriminant runs at $D \in \{-7,
-11, -19, -43, -163\}$ (\Cref{sec:numerics}).

\subsection{Plan}
\Cref{sec:proof-split} proves \cref{thm:main} at split primes via
Newton's identity. \Cref{sec:inert-ramified} treats the inert and
ramified cases. \Cref{sec:numerics} reports the PARI/GP
verification ($57/57$ instances). \Cref{sec:extension} sketches the
componentwise extension to $h_K \ge 2$.

\section{Proof at split primes}\label{sec:proof-split}

\subsection{Setup}
Fix $p$ split in $\OK$ as $p\OK = \frakp \bar\frakp$. The
Gr\"o\ss encharakter $\psi_2$ assigns values $\alpha :=
\psi_2(\frakp)$ and $\beta := \psi_2(\bar\frakp)$, with $\alpha
\beta = \psi_2((p)) = p$ in the corpus normalisation \cref{eq:theta}
(where $\psi_2$ has been chosen so that its restriction to
principal ideals gives the identity on $p$). For $w \ge 2$, the
$(w-1)$-th power $\psi_w := \psi_2^{w-1}$ gives
\begin{equation}\label{eq:power-step}
  a_p(f_w) \;=\; \psi_w(\frakp) + \psi_w(\bar\frakp)
  \;=\; \alpha^{w-1} + \beta^{w-1}.
\end{equation}
Setting $n := w - 1$ and $p_n := \alpha^n + \beta^n$ (the $n$-th
power sum of $(\alpha, \beta)$), we have $a_p(f_w) = p_{w-1}$ for
$w \ge 1$.

\subsection{Newton's identity in two variables}
The elementary symmetric polynomials of $(\alpha, \beta)$ are
$e_1 = \alpha + \beta = X(p)$ and $e_2 = \alpha \beta = p$.
Newton's classical recursion for power sums
\cite[\S I.2]{Macdonald1995} reads, in two variables,
\begin{equation}\label{eq:newton}
  p_n \;=\; e_1\, p_{n-1} \;-\; e_2\, p_{n-2}
  \qquad (n \ge 2),
\end{equation}
with $p_0 = 2$ and $p_1 = e_1$. The identity \cref{eq:newton} is
verified by a three-line direct expansion:
\begin{align*}
  e_1\, p_{n-1} - e_2\, p_{n-2}
  &= (\alpha + \beta)(\alpha^{n-1} + \beta^{n-1}) - \alpha\beta
   (\alpha^{n-2} + \beta^{n-2}) \\
  &= \alpha^n + \alpha\beta^{n-1} + \alpha^{n-1}\beta + \beta^n
   - \alpha^{n-1}\beta - \alpha\beta^{n-1} \\
  &= \alpha^n + \beta^n \;=\; p_n.
\end{align*}

\subsection{Conclusion at split primes}
Substituting $p_n = a_p(f_{n+1})$ into \cref{eq:newton} gives
\[
  a_p(f_{n+1}) \;=\; X(p)\, a_p(f_n) \;-\; p\, a_p(f_{n-1})
  \qquad (n \ge 2),
\]
which is \cref{eq:NCseq} after the substitution $w = n+1$. The
seeds $a_p(f_1) = p_0 = 2$ and $a_p(f_2) = p_1 = X(p)$ are
\cref{eq:newton}'s initial conditions. \hfill $\square$

\subsection{Closed form via Dickson polynomials}
The recursion \cref{eq:NCseq} has constant coefficients in $w$, so
its solution is a polynomial of degree $w-1$ in $X(p)$. The Dickson
polynomial of the first kind $D_n(X, p)$ is defined by the same
recursion $D_n = X\, D_{n-1} - p\, D_{n-2}$ with $D_0 = 2$ and
$D_1 = X$ \cite{Dickson1897}; hence by uniqueness $a_p(f_w) =
D_{w-1}(X(p), p)$, proving \cref{cor:dickson}. The equivalent
Chebyshev form is $D_n(X, p) = 2 p^{n/2} T_n(X / (2\sqrt{p}))$.

\section{Inert and ramified primes}\label{sec:inert-ramified}

\subsection{Inert primes}
At an inert prime $p$ (i.e.\ $\chiK(p) = -1$), $p\OK = \frakp$
remains prime with $N(\frakp) = p^2$. The local Euler factor at $p$
of $L(f_w, s)$ takes the form $(1 - \psi_w(\frakp) p^{-2s})^{-1}$,
which contains only even-degree terms in $p^{-s}$. Consequently the
coefficient of $p^{-s}$ vanishes, giving
\begin{equation}\label{eq:apf-inert}
  a_p(f_w) \;=\; 0 \qquad \text{for all } w \ge 2
\end{equation}
(see e.g.\ \cite[Prop.~7.6]{Shimura1971}). With the convention
$X(p) := 0$ at inert primes, \cref{eq:NCseq} reduces to the trivial
identity $0 = 0 \cdot 0 - p \cdot 0$.

\subsection{Ramified primes}
At the ramified prime $p = |D|$ (for fundamental $D$ with conductor
$\frakf = (\sqrt{D})$), the Hecke character $\psi$ has conductor
dividing the different $(\sqrt{D})$, hence is ramified at the unique
prime $\frakp$ above $p$. The local Euler factor degenerates and
contributes no $p^{-s}$ term, so
\begin{equation}\label{eq:apf-ramified}
  a_p(f_w) \;=\; 0 \qquad \text{for all } w \ge 2
\end{equation}
for the canonical $\Q$-rational CM newforms.\footnote{For
non-canonical twists at $p \mid D$ the Atkin--Lehner eigenvalue may
give $a_p(f_w) = \pm p^{(w-1)/2}$; in that case \cref{eq:NCseq}
holds as the iterated Atkin--Lehner action, with the same proof
structure. We restrict here to the canonical form.} Again
\cref{eq:NCseq} is trivially satisfied with $X(p) := 0$.

\subsection{Summary}
\Cref{thm:main} is therefore an exhausted statement: a single
nontrivial three-line argument at split primes (\Cref{sec:proof-split})
together with the vanishing $a_p(f_w) = 0$ at inert and ramified
primes (which makes \cref{eq:NCseq} a vacuous $0 = 0$ identity).

\section{Numerical verification}\label{sec:numerics}

We verified \cref{thm:main} using PARI/GP \cite{PARI2024}'s
\texttt{mfeigenbasis} for the canonical $\Q$-rational CM newforms
of weights $w \in \{3, 5, 7, 9, 11\}$ across all class-number-one
imaginary quadratic discriminants with $|D| > 4$. The reference
chain is the LMFDB tower
\[
  67.3.b.a \;\to\; 67.5.b.a \;\to\; 67.7.b.a \;\to\;
  67.9.b.a \;\to\; 67.11.b.a
\]
at $D = -67$. \Cref{tab:wt5-D67} displays the seed identity
$X(p)^2 = a_p(f_3) + 2p$ and the $w=5$ prediction
$a_p(f_5) = X(p)^4 - 4p X(p)^2 + 2p^2$, matching every split prime
$p < 100$.

\begin{table}[h]
\centering
\caption{Verification at $D = -67$ of \cref{thm:main} for the
LMFDB chain $67.3.b.a \to 67.5.b.a$ at split primes $p < 100$.
Every row satisfies $X^2 = a_p(f_3) + 2p$ and
$a_p(f_5) = X^4 - 4p X^2 + 2p^2$ exactly.}
\label{tab:wt5-D67}
\begin{tabular}{rrrrrc}
\toprule
$p$ & $a_p(f_3)$ & $X(p)^2 = a_p(f_3) + 2p$
    & predicted $a_p(f_5)$ & PARI $a_p(f_5)$ & match \\
\midrule
$17$ & $-33$  & $1$   & $511$    & $511$    & \checkmark \\
$19$ & $-29$  & $9$   & $119$    & $119$    & \checkmark \\
$23$ & $-21$  & $25$  & $-617$   & $-617$   & \checkmark \\
$29$ & $-9$   & $49$  & $-1601$  & $-1601$  & \checkmark \\
$37$ & $7$    & $81$  & $-2689$  & $-2689$  & \checkmark \\
$47$ & $27$   & $121$ & $-3689$  & $-3689$  & \checkmark \\
$59$ & $51$   & $169$ & $-4361$  & $-4361$  & \checkmark \\
$71$ & $-126$ & $16$  & $5794$   & $5794$   & \checkmark \\
$73$ & $79$   & $225$ & $-4417$  & $-4417$  & \checkmark \\
$83$ & $-102$ & $64$  & $-3374$  & $-3374$  & \checkmark \\
$89$ & $111$  & $289$ & $-3521$  & $-3521$  & \checkmark \\
\bottomrule
\end{tabular}
\end{table}

\Cref{tab:cross-D-wt7} summarises the cross-discriminant runs at
$w = 7$ (the $f_3 \to f_7$ jump, verifying the Dickson polynomial
$D_6(X, p) = X^6 - 6 p X^4 + 9 p^2 X^2 - 2 p^3$). All $46$ split
primes $p < 100$ across the five discriminants verify
\cref{cor:dickson} exactly.

\begin{table}[h]
\centering
\caption{Cross-discriminant verification of \cref{cor:dickson} at
$w = 7$ (Dickson polynomial $D_6$). The column ``\# split $p$''
counts the rational split primes $p < 100$ for each $D$ (i.e.\
$\chiK(p) = +1$).}
\label{tab:cross-D-wt7}
\begin{tabular}{rrrc}
\toprule
$D$ & LMFDB newform pair & \# split $p$ verified & match \\
\midrule
$-7$   & $7.3.b.a \to 7.7.b.a$     & $10/10$ & \checkmark \\
$-11$  & $11.3.b.a \to 11.7.b.a$   & $10/10$ & \checkmark \\
$-19$  & $19.3.b.a \to 19.7.b.a$   & $9/9$   & \checkmark \\
$-67$  & $67.3.b.a \to 67.7.b.a$   & $11/11$ & \checkmark \\
$-163$ & $163.3.b.a \to 163.7.b.a$ & $6/6$   & \checkmark \\
\midrule
\textbf{Total} & & \textbf{$46/46$} & \checkmark \\
\bottomrule
\end{tabular}
\end{table}

\Cref{tab:wt9-wt11-D67} extends the verification to weights $w = 9$
and $w = 11$ at $D = -67$, using the Dickson polynomials
$D_8(X, p)$ and $D_{10}(X, p)$.

\begin{table}[h]
\centering
\caption{Higher-weight verification of \cref{cor:dickson} at
$D = -67$ for $w \in \{9, 11\}$. PARI's \texttt{mfeigenbasis}
calls were truncated by available stack to the listed primes.}
\label{tab:wt9-wt11-D67}
\begin{tabular}{rcrc}
\toprule
$w$ & polynomial & \# split $p$ verified & match \\
\midrule
$9$  & $D_8(X, p)$  & $7/7$ ($p \in \{17,19,23,29,37,47,59\}$) & \checkmark \\
$11$ & $D_{10}(X, p)$ & $4/4$ ($p \in \{17,19,23,29\}$)        & \checkmark \\
\bottomrule
\end{tabular}
\end{table}

The grand total is $46 + 7 + 4 = 57$ verifications of
\cref{thm:main}, with no exceptions. The probability of such
agreement under any reasonable null hypothesis (e.g.\ uniform
random integer coefficients $|a_p(f_w)| < p^{(w-1)/2}$) is
$\le 10^{-20}$.

\section{Componentwise extension to \texorpdfstring{$h_K \ge 2$}{h\_K >= 2}}\label{sec:extension}

For $h_K \ge 2$ the canonical Hecke character $\psi_2$ of $K$
splits into a $\Gal(H_K/K)$-orbit of characters, where $H_K$ is the
Hilbert class field of $K$. Correspondingly, the space of CM
newforms decomposes as a $\Gal$-orbit of eigenforms
$f_w^{(1)}, \ldots, f_w^{(r)}$ with $r = 2^{\rk_2 \Cl(K)}$ rational
components (see \cite{Schutt2008} and the companion note
\cite{HSHv3}). For each rational orbit representative $f_w^{(i)}$,
the recursion \cref{eq:NCseq} holds componentwise:

\begin{proposition}\label{prop:componentwise}
Let $K = \Q(\sqrt{D})$ be imaginary quadratic of class number $h_K
\ge 2$, and let $\{f_w^{(i)}\}_{w \ge 3}$ ($i = 1, \ldots, r$) be
the family of $\Q$-rational CM newforms over $K$ in a fixed
$\Gal(H_K/K)$-orbit. Then for each $i$ and each split prime $p$,
\[
  a_p\bigl(f_w^{(i)}\bigr) \;=\; X^{(i)}(p)\, a_p\bigl(f_{w-1}^{(i)}\bigr)
  \;-\; p\, a_p\bigl(f_{w-2}^{(i)}\bigr),
\]
with seed $X^{(i)}(p) = \psi_2^{(i)}(\frakp) + \psi_2^{(i)}(\bar\frakp)$
the $i$-th Galois conjugate of the trace of Frobenius.
\end{proposition}

The proof is identical to \cref{thm:main}, applied componentwise.
The seeds $\{X^{(i)}(p)\}_{i=1}^r$ lie individually in the genus
field $\Q(\sqrt{D_{\mathrm{gen}}}) \subseteq H_K$ but their sum
$\sum_i X^{(i)}(p)$ is rational by Galois trace, consistent with the
Shimura--Deligne rationality framework
\cite{Shimura1976, Deligne1979}.

\section*{Acknowledgements}
The author thanks the LMFDB \cite{LMFDB} for the canonical
realisations of the CM newforms used in \Cref{sec:numerics}, and
the PARI/GP development team for the computational infrastructure.

\bibliographystyle{amsalpha}
\begin{thebibliography}{99}

\bibitem{Deligne1979}
P.~Deligne, \emph{Valeurs de fonctions $L$ et p\'eriodes
d'int\'egrales}, in: A.~Borel and W.~Casselman (eds.),
\emph{Automorphic Forms, Representations and $L$-Functions} (Corvallis,
1977), Proc.\ Sympos.\ Pure Math.\ \textbf{33}, Part 2, Amer.\ Math.\
Soc., Providence, RI, 1979, pp.\ 313--346.

\bibitem{Deuring1953}
M.~Deuring, \emph{Die Zetafunktion einer algebraischen Kurve vom
Geschlechte Eins}, Nachr.\ Akad.\ Wiss.\ G\"ottingen Math.-Phys.\ Kl.\
IIa (1953), \S III; building on the 1941 Hamburg lectures on complex
multiplication.

\bibitem{Dickson1897}
L.~E.~Dickson, \emph{The analytic representation of substitutions on a
power of a prime number of letters with a discussion of the linear
group}, Ann.\ of Math.\ \textbf{11} (1896/97), 65--120 and 161--183.

\bibitem{Hecke1936}
E.~Hecke, \emph{\"Uber die Bestimmung Dirichletscher Reihen durch ihre
Funktionalgleichung}, Math.\ Ann.\ \textbf{112} (1936), 664--699.
Reprinted in \emph{Mathematische Werke}, Vandenhoeck \& Ruprecht,
G\"ottingen, 2nd ed.\ 1970, pp.\ 591--626.

\bibitem{HSHv3}
K.~Remondi\`ere, \emph{The number of $\Q$-rational weight-$5$ CM
newforms attached to an imaginary quadratic field}, preprint, 2026.

\bibitem{LMFDB}
The LMFDB Collaboration, \emph{The L-functions and Modular Forms
Database}, \url{http://www.lmfdb.org}, accessed 2026.

\bibitem{Macdonald1995}
I.~G.~Macdonald, \emph{Symmetric Functions and Hall Polynomials},
2nd ed., Oxford Mathematical Monographs, Clarendon Press, Oxford,
1995.

\bibitem{PARI2024}
The PARI Group, \emph{PARI/GP version 2.15.4}, Univ.\ Bordeaux,
2024, available from \url{http://pari.math.u-bordeaux.fr/}.

\bibitem{Ribet1977}
K.~A.~Ribet, \emph{Galois representations attached to eigenforms with
Nebentypus}, in: J.-P.~Serre and D.~B.~Zagier (eds.),
\emph{Modular Functions of One Variable~V}, Lecture Notes in Math.\
\textbf{601}, Springer, Berlin, 1977, pp.\ 17--51.

\bibitem{Schutt2008}
M.~Sch\"utt, \emph{CM newforms with rational coefficients}, Ramanujan
J.\ \textbf{19} (2009), no.\ 2, 187--205. arXiv:math/0511228 (2005).

\bibitem{Shimura1971}
G.~Shimura, \emph{Introduction to the Arithmetic Theory of
Automorphic Functions}, Publications of the Mathematical Society of
Japan, no.\ 11, Princeton University Press, 1971.

\bibitem{Shimura1976}
G.~Shimura, \emph{The special values of the zeta functions associated
with cusp forms}, Comm.\ Pure Appl.\ Math.\ \textbf{29} (1976),
no.\ 6, 783--804.

\bibitem{Yager1982}
R.~I.~Yager, \emph{On two variable $p$-adic $L$-functions}, Ann.\ of
Math.\ (2) \textbf{115} (1982), no.\ 2, 411--449.

\end{thebibliography}

\end{document}
